% 本文件由 LaTeX 排版 Agent 根据有效 translation.json 自动生成。
% author=Commodianus; work_id=0411; title=论基督徒纪律

\chapter{论基督徒纪律}

% structure: p0001 source=h1 target=omitted reason=page-title-already-rendered-by-parent

\Emph{（以藏头诗形式表达）}\par

约公元二四〇年写成。我们的作者似乎是一位北非主教，关于他除了从他本人的著作中得知的以外，所知甚少。有人推测他倾向于\Person{普拉克塞亚斯}的某些观点，也倾向于千禧年论者，但或许理由不充分。他的千禧年主义反映了非常早期时代的观点，而没有后来时期的腐败千禧年主义——后者导致了对整个体系的实际否定。在他的著作中，只留存下来两首诗，其中第二首是最近发现的，未收入\WorkTitle{爱丁堡系列}。我深感遗憾，它未能收入我们的系列中。\par

作为一部诗作，以下散文译本或许并未对其不公。他的诗律被认为非常晦涩，其用词是北非低劣的\OriginalTerm{方言}{patois}。但在这片古代的残篇中，一位实践派基督徒的虔诚与热忱似乎处处可见。\par

% structure: p0005 source=h2 target=section reason=relative-heading-rank; base=h2

\section{I. 前言。}

我的前言向迷途者指明道路，当生命的终点到来时给予善意的眷顾，使他得以成为永恒——这是无知的心所不信的。我也曾长久迷失，因追随\OriginalTerm{异教}{heathen}神庙，我父母自己也蒙昧无知。后来藉着研读关于法律的书，我才从中脱离。我为主作证；唉，我为众公民的群体悲伤！他们不知道自己在寻求虚无的神明时失去了什么。藉这些事深受教导，我以此真理教导无知者。\par

% structure: p0007 source=h2 target=section reason=relative-heading-rank; base=h2

\section{II. 天主的义怒。}

在法律中，天地海的主曾命令说：不可敬拜你们亲手用木头或金子所造的虚无神明，免得我的愤怒因这些事毁灭你们。在\Person{梅瑟}之前的百姓，不谙法律，无法度，也不认识天主，向毁灭的神明祈祷，按照它们的模样制造虚无的偶像。主将犹太人从埃及地领出后，随即给他们颁赐了法律；全能者吩咐这些事，要他们惟独事奉祂，不事奉那些偶像。此外，在那法律中也教导了关于复活、以及重新在世上幸福生活的盼望——只要离弃虚无的偶像，不敬拜它们。\par

% structure: p0009 source=h2 target=section reason=relative-heading-rank; base=h2

\section{III. 对魔鬼的敬拜。}

当全能的天主为美化世界的本性，愿意那大地被天使眷顾时，他们被派下来却藐视祂的法律。妇女的美貌使他们偏离；因此，他们被玷污，不能再返回天堂。他们背叛天主，口出恶言反对祂。于是至高者宣告了对他们的审判；据说从他们的后裔中生出了巨人。藉着他们，技艺在地上被显明，他们教授染羊毛以及一切所作的事；当他们死去时，人为他们竖立了偶像。但全能者因为他们属于邪恶的种子，不允许他们死后从死亡中返回。因此，如今他们游荡，败坏许多身体；尤其是这些，你们今日敬拜并祈求为神。\par

% structure: p0011 source=h2 target=section reason=relative-heading-rank; base=h2

\section{IV. 萨图尔努斯。}

老\Person{萨图尔努斯}，如果他是神，怎么变老？或者如果他是神，为何因恐惧而吞噬自己的儿女？但因为他不是神，他疯狂地吞食了儿子们的内脏。他曾是地上的王，生在\Place{奥林波斯山}上；他并非神圣，却自称为神。他陷入心智软弱，吞下一块石头代替儿子。于是他成了神；后来被称为\Person{朱庇特}。\par

% structure: p0013 source=h2 target=section reason=relative-heading-rank; base=h2

\section{V. 朱庇特。}

这位朱庇特生在\Place{布里塔岛}，为萨图尔努斯所生；长大后，他剥夺了父亲的王位。他诱骗了贵族们的妻子和姊妹。此外，铁匠\Person{皮拉克蒙}为他打造了一根权杖。起初，天主创造了天、地、海。但那可怕的东西，生于时间之中，从洞穴中走出为青年，被秘密喂养。看，那位天主是万有的创造者，不是那朱庇特。\par

% structure: p0015 source=h2 target=section reason=relative-heading-rank; base=h2

\section{VI. 论同一位朱庇特的雷霆。}

你们说，愚人啊，朱庇特打雷。是他投掷雷霆；如果幼稚的心这样想，为何你们两百年仍是婴孩？并且还要永远这样？婴孩期已进入成熟，老年不喜爱琐事，童年已逝；少年的心也当离去。你们的思想应当属于成人的品格。那么你们是愚人，竟相信朱庇特打雷。他生在地上，用山羊奶喂养。因此，如果萨图尔努斯吞了他，那么他死后，那些时候谁降雨呢？尤其，如果一位神被认为生于有死的父亲，那么萨图尔努斯在地上衰老，在地上死。没有人预言他先前的出生。或者如果他打雷，法律就应由他颁布。诗人们虚构的故事引诱了你们。然而，他在\Place{克里特}统治，并死在那里。你们视为全能者的他，曾成为\Person{阿尔克墨涅}的情人；如果他活着，如今也照样会爱上活人。你们向不洁的神明祈祷，称那些从巨人那里由有死的种子所生的为天上的神。你们听见并读到，他生在地上：那个败坏者凭什么那么配升天？而且独眼巨人据说为他锻造了雷霆；尽管他不死，却从有死者那里接受武器。你们凭自己的权威，将犯下如此多罪行、并且还是杀害亲人的凶手的人，送上了天堂。\par

% structure: p0017 source=h2 target=section reason=relative-heading-rank; base=h2

\section{VII. 论\OriginalTerm{七层神坛}{Septizonium}与星辰。}

你缺乏智慧，致使你被天区之圈子所欺骗，也许由此你发现你必须祈求朱庇特。那里也提到萨图尔努斯，但只是作为一颗星，因为他是被朱庇特驱逐的，或者让朱庇特被相信存在于那颗星中。那位控制极地星座者、土地的播种者；那位与特洛伊人作战者，他爱美丽的维纳斯。或在众星之间，玛尔斯因婚姻的嫉妒与她一同被捉：他被称为青年之神。哦，极其愚蠢，竟以为那些由迈亚所生者从星辰统治世界，或他们统治世界的整个本性！他们自身受制于创伤，生活在命运的支配之下，淫秽、好奇、过着不敬虔生活的战士；并且他们生了儿子，与自己一样会死，且全都可怕、愚蠢、强壮，在七重环带中。如果你崇拜星辰，也崇拜十二个\Emph{黄道带}星座，以及白羊座、金牛座、双子座，还有凶猛的狮子座；最后，它们延伸成双鱼座——煮熟它们，你就会证明它们。一条没有律法的律法是你的避难所：想要成为的，就会占上风。一个女人想要放荡；她寻求无拘无束地生活。你们自己将成为你们所希望成为的，并祈求为神和女神。我过去在迷途中就是这样崇拜，如今我谴责它。\par

% structure: p0019 source=h2 target=section reason=relative-heading-rank; base=h2

\section{八、论太阳和月亮。}

论到太阳和月亮，你们犯了错误，尽管它们就在我们眼前；因为你们，像我从前一样，认为必须向它们祈祷。它们确实是在星辰之中；但它们并非自行运行。全能者在起初创造万物时，在第四天把它们与星辰一同安置在那里。而且，他确实在律法中命令任何人都不可崇拜它们。你们崇拜如此多的神，他们对于生命毫无应许，他们的律法不在大地上，他们本身也未被预言。但少数几位司祭误导你们，他们说任何注定要死的神祇都可以有用。现在靠近，阅读，学习真理。\par

% structure: p0021 source=h2 target=section reason=relative-heading-rank; base=h2

\section{九、墨丘利。}

让你们的墨丘利被描绘成穿着萨拉巴卢姆，头盔或帽子上有翅膀，其余部分赤裸。我看见一件奇妙的事，一位神祇带着一个小袋子飞翔。跑吧，可怜的人们，张开你们的衣兜，当他飞过时，让他倒空他的袋子：从那时起就预备好。看那画中的神，因为他会这样从高处向你们抛掷金钱：然后安然起舞。虚妄的人啊，你岂不疯狂，竟崇拜天上画出来的神祇？你若不知如何生活，就继续与野兽同居吧。\par

% structure: p0023 source=h2 target=section reason=relative-heading-rank; base=h2

\section{十、尼普顿。}

你们使尼普顿成为一位源自萨图尔努斯的神；他手持三叉戟，以便刺鱼。从他这样的装备来看，显然他是海神。难道不是他自己与阿波罗一起为特洛伊人修筑城墙吗？那个可怜的石匠如何成为神？难道他没有生出独眼巨人怪物？他自己死后难道不能复活，尽管他的构造允许这样？既然这样被生，他便生了已经死过一次的。\par

% structure: p0025 source=h2 target=section reason=relative-heading-rank; base=h2

\section{十一、阿波罗——占卜者与假神。}

你们使阿波罗成为一位西塔拉琴演奏者和占卜者。他最初生于迈亚，在提洛岛；随后，为得工钱，作为建造者，听命于拉俄墨冬王，修筑了特洛伊人的城墙。他自立为神，而你们被引诱，以为他是神；在他骨子里，卡珊德拉的爱情燃烧着，那处女狡猾地戏弄了他；尽管是神，他却受骗。凭他占卜官的职务，他本应识破那心怀二意者。此外，被拒绝以后，他虽是神，却从那里离开。那处女以其美貌使他燃烧，他本应烧死她；而她本应首先爱那如此好色地开始爱上达芙妮的神，他仍追赶她，想要强暴那少女。那愚人徒然恋爱。他无法通过奔跑得到她。诚然，他若是神，就会穿越空气追上她。她先到了屋檐下，而那位神留在外面。人的种族欺骗你们，因为他们生活方式可悲。此外，据说他曾放牧阿德墨托斯的牛群。在规定的游戏中，他将铁饼抛向空中，却无法在它落下时制住它，它砸死了他的朋友。那是他同伴伊阿辛托斯的末日。他若是神，本应预知朋友的死亡。\par

% structure: p0027 source=h2 target=section reason=relative-heading-rank; base=h2

\section{十二、利伯尔神——巴克斯。}

你们自己说利伯尔神确实是两次被生。首先，他生在印度，由普洛塞庇娜和朱庇特所生；与泰坦作战时，他的血流出，他就像凡人一样断气了。再次，从他死亡中复原，在另一个子宫里，塞墨勒又由朱庇特怀了他，这是第二个迈亚；她的子宫被剖开，他从已死的母亲那里被取出，接近出生，作为一个婴孩交给尼索斯抚养。因这双重出生，他被称作狄俄尼索斯；他的宗教被虚妄地错误遵守；他们举行他的秘仪，以致如今他们自己似乎要么是鲁莽的，要么是米姆涅尔莫墨鲁斯的戏仿者。他们同谋作恶；他们预先演练假装的狂热，以便欺骗别人说一位神祇在场。由此你们清楚地看见，人们过着像他一样的生活，被他亲手压榨的酒猛烈激动；他们在醉酒纵欲中赐予他神圣的尊荣。\par

% structure: p0029 source=h2 target=section reason=relative-heading-rank; base=h2

\section{十三、不可战胜者。}

不可战胜者是从岩石中诞生的，如果他被视为神的话。那么，现在告诉我们，反之，这两者中哪个在先？岩石胜过了神：那么就必须寻找岩石的创造者。此外，你们还把他描绘成一个贼；尽管，他若是神，当然不靠偷窃为生。他确凿无疑属于大地，具有怪物的本性。他把别人的牛带进他的洞穴；正如武尔坎的儿子卡库斯所做的那样。\par

% structure: p0031 source=h2 target=section reason=relative-heading-rank; base=h2

\section{十四、西尔瓦努斯。}

再者，西尔瓦努斯从何看起来是神呢？也许从这一点来\Emph{如此称呼他}是合宜的，因为笛子甜美地歌唱，因为他赐予木材；因为，也许并非如此。你们买了一个可买的主人，当你们从他那里买了之后。看哪，木材没有了！他该得什么？愚人啊，你们不羞于崇拜这样的画像吗？寻求一位独一天主，他会允许你们死后生存。离开那些在生命中已成为死者吧。\par

% structure: p0033 source=h2 target=section reason=relative-heading-rank; base=h2

\section{十五、赫拉克勒斯。}

\Person{赫拉克勒斯}因摧毁了惯于偷窃\Person{埃万德}牛群的\Place{阿文廷山}上的怪物，\Emph{被视为神}：未受教化的乡野人心，当他们想要感谢而非赞美那位不在场的雷神时，愚昧地向一位应祈求的神献上牺牲作为还愿，他们为自己立下了乳香的祭台作为纪念。由此而来，他以古老的方式受到敬拜。但他并非神，尽管他武力强大。\par

% structure: p0035 source=h2 target=section reason=relative-heading-rank; base=h2

\section{第十六章 论诸神与女神}

你们说那些显然是残忍的神是神，又说星宿决定你们的命运。那么，请说首先向谁举行祭祀。在道路两旁，夭折的死亡在游荡。如果星宿赋予生命，你们为何向神祈祷？你徒然受骗，寻求恳求亡灵，并称那些被造之物为主宰。或者，我不知道你们向哪些女子祈祷为女神——\Person{贝娄娜}、\Person{涅墨西斯}，连同天界复仇女神、贞女和\Person{维纳斯}，你们的妻子因她们而腰腿软弱。此外，巷中还有其他未被计数的邪灵，他们戴在颈上，以致他们自己也无法说明。瘟疫倒应被驱逐到地极。\par

% structure: p0037 source=h2 target=section reason=relative-heading-rank; base=h2

\section{第十七章 论他们的像}

少数邪恶而空洞的诗人欺骗你们；他们艰难谋生时，以神秘的伪装为他人装饰谎言。于是假装被某神击打，他们歌颂祂的威严，并劳累于祂的形体之下。你们常看见丁底摩利亚祭司，他们喧嚷着进入奢靡，假装疯狂，或者用污秽的斧头击打自己的背，尽管他们的教导使他们用血来医治。看啊，他们以何名不强迫那些以健全心智首先与自己联合的人。但为取走礼物，他们寻求这样的心智。于是看一切如何都是伪造的。他们给单纯的人民投下阴影，以免他们相信——当他们在虚空中灭亡时——从古代流传下来的一件事，好让说假话的先知被相信；但他们的威严一言不发。\par

% structure: p0039 source=h2 target=section reason=relative-heading-rank; base=h2

\section{第十八章 论阿米达特斯与大神}

我们已经说了许多关于可憎迷信的事，但仍继续探讨，以免被认为遗漏了什么。敬拜者按自己的方式敬拜他们的阿米达特斯。当庙中有金子时，他对他们来说是伟大的。他们将自己的头置于他的权下，仿佛他亲临。事情到了极点，\Person{凯撒}取走了金子。神祇衰败、逃走或化为火焰。这邪恶的始作俑者昭然若揭，他造了这位神，虚假预言误导了众多伟人，却只对那习惯为神者沉默。因为声音迸发，仿佛变了心，好像木神在对他耳语。现在你们自己说，难道它们不是假神吗？从那个怪异中，那位先知毁灭了多少人？他忘了预言，而他从前习惯预言；所以那些怪物在贪杯者中伪造，他们可恶的狂妄造出神祇，因为那些像被抬来抬去，干枯了。因为他自己沉默，也无人预言关于他任何事。但你们甘愿自取灭亡。\par

% structure: p0041 source=h2 target=section reason=relative-heading-rank; base=h2

\section{第十九章 论虚妄的涅墨西亚奇}

这难道不是耻辱吗？明智的人被误导，敬拜这样一个神，或者说一块木头是\Person{狄安娜}？你相信一个早晨醉醺醺、便秘、濒死的人，他工巧地编造所见。当他严厉生活时，他吞噬自己的内脏。一个可憎者玷污所有公民；他聚集相似的一群，与他一同伪造故事，以美化一位神。他不知如何为自己预言；却敢为他人预言。他高兴时把它扛在肩上，又放下来。他旋转，被自己的双叉树木转动，仿佛你会认为他被树木的神明感动。你们不敬拜他们自己假报的神；你们敬拜祭司本人，徒然惧怕他们。但若你们心坚，就立即逃离死亡的庙宇。\par

% structure: p0043 source=h2 target=section reason=relative-heading-rank; base=h2

\section{第二十章 提坦神族}

你们称提坦为你们的\OriginalTerm{图坦斯}{Tutans}。你们请求这些凶悍者在你们屋顶下沉默，如同众多拉瑞斯、神龛、提坦式的像。因为你们愚昧地敬拜那些死于恶死的人，却不读自己的律法。他们自己不说话，你们却敢称那些从铜器熔铸的为神；你们倒不如将它们熔铸成小器皿给自己。\par

% structure: p0045 source=h2 target=section reason=relative-heading-rank; base=h2

\section{第二十一章 蒙泰西亚尼}

你们也称山为神。让它们在金中统治，被恶玷污，以回避之心相助。因为若你们仍有纯洁的精神和平静的心思，你们本应自行审察它们。你们作为人已变得无知，若以为这些能救你们，无论它们统治与否。若你们寻求任何健全之物，反要求律法的义，它带来救恩的帮助，并说你们将成永恒。因为你们在虚空中追随之事，暂时令你们欢喜。你们欢乐片刻，随后在深渊中哀哭。离开这些吧，若你们愿与基督一同复活。\par

% structure: p0047 source=h2 target=section reason=relative-heading-rank; base=h2

\section{第二十二章 时代的愚钝}

哀哉，我悲痛，公民们，你们竟被此世如此蒙蔽。有人奔赴抽签，有人凝视飞鸟，有人宰杀咩咩牲畜，流其鲜血，召请亡灵，并轻信地渴望听闻虚妄的回应。当如此多的领袖和君王就生命之事商议时，知晓其预兆对他们有何益处？我恳请你们，公民们，学习何为善；当心偶像神庙。诚然，你们众人当在全能者的法律中寻求。主的主在诸天中如此乐意，让魔鬼为我们的操练而游荡于世。然而，另一方面，祂已发出祂的诫命：凡离弃其祭台的，将成为天堂的居民。因此，我不必在一篇短论中详加辩论。法律教导；它在你们当中呼唤你们。你们自己思量。你们已走上了两条路；当抉择那正确的。\par

% structure: p0049 source=h2 target=section reason=relative-heading-rank; base=h2

\section{那些随时备妥的人}

当你顺从肚腹时，你声称自己无罪；并且，仿佛彬彬有礼，使自己随时准备就绪。祸哉，愚昧的人！你亲自在死亡四周观望。你以一种野蛮的方式寻求无律法而活。你甚至自诩玩弄言辞，装作单纯。我与这样一个人过着纯朴的生活。你相信自己在活着，而其实你只想填满肚腹。在你家中不光彩地坐下，毫无价值，随时准备宴饮，却逃避诫命。或者因为你相信天主不会审判死者，你便愚昧地自命为天堂的统治者，取代祂。你视你的肚腹仿佛你能供养它。你一时看似亵渎，一时又似圣洁。你以暴君的面貌，显现为天主的恳求者。你将在你的命运中感受到，你由谁的法律所扶助。\par

% structure: p0051 source=h2 target=section reason=relative-heading-rank; base=h2

\section{那些活在两者之间的人}

你以模棱两可的方式活在两者之间，自以为有所提防，继续前行，脱去了法律，被奢侈击垮。你徒然期待如此多的事物，为何寻求不义之事？凡你所行的，死后将留在那里。愚昧的人啊，你思量：你昔日并不存在，看，如今你显现了。你不知从何而来，也不知从何得滋养。你回避你生命之优良与仁慈的天主，以及你的治理者，祂更愿你存活。你转向自己，将背转向天主。你将自己沉入黑暗，却自以为在光明中存留。你为何跑到会堂里的法利塞人那里，好使祂怜悯你，而你却自愿否认祂？然后你再次外出；你寻求有益之事。你愿活在两条路之间，但由此你将灭亡。并且你还说：\Quote{那从死亡中救赎者是谁，使我们相信祂，既然那里有刑罚？}啊！恶毒的人啊，不会如你所想的。因为生活得好的人，死后有益。然而，当你一日死去，你将被带到邪恶的地方。但相信基督的人，将被引入良善的地方，那些得了此喜悦的人得蒙抚慰；但对你，心怀二意的人，惩罚是不在身上的。折磨者的驱策激起你向你弟兄喊叫。\par

% structure: p0053 source=h2 target=section reason=relative-heading-rank; base=h2

\section{那些恐惧而不肯相信的人}

愚昧的人啊，你何时才承认基督？你避开肥沃的田地，将种子撒在不毛之地。你寻求停留在盗贼盘桓的树林中。你说：\Quote{我也是属于天主的}；却浪迹户外。如今，在这么多次邀请之后，进入宫殿吧。现今庄稼已熟，时节屡次预备了。看，现在收割吧！什么？你还不悔改？那么现在，若你还没有，就收取合时的葡萄。相信而得生命的时刻，正在死亡的时刻中。天主的第一条法律是后来法律的基础。它指派你相信第二条法律。威胁并非来自祂本身，而是从这法律来辖制你。如今震惊了吧，起誓说你将相信基督；因为《旧约》论及祂而宣告。因为只需要相信那已死的人，才能复活而永远活着。因此，若你不信这些事将要成就，最终他将在第二次死亡中为他的罪所胜。我将在这一小篇论著中宣告将来的事。从中可知何时当优先怀抱希望。我仍劝勉你尽快相信基督。\par

% structure: p0055 source=h2 target=section reason=relative-heading-rank; base=h2

\section{致那些抗拒生活天主之基督的法律的人}

不幸的人哪，你拒绝了天上的纪律，却想不受约束地游荡，冲向死亡。奢侈和世上短暂的快乐正在毁灭你，因此你将在地狱中永远受苦。你愚蠢地陶醉于那些虚妄的快乐。难道这些不能使你成为死人吗？难道三十年的光阴终究不能使你成为明智的人吗？你不知道自己最初如何迷失，看看上古的时间，你现在以为可以在罪恶中享受快乐的生活。这些是你朋友的毁灭、战争、邪恶的欺诈、带血的盗窃：身体被疮痍折磨，呻吟和哀嚎不断；无论是小病侵袭你，还是久病缠身，或是丧子，或是为失去的妻子哀哭。一切都是荒凉：哀哉，尊荣因恶习和贫穷而从高位坠落；如果你久久衰弱，这无疑是加倍的。当这如玻璃般易碎的生命终有一死时，你还称之为生命吗？现在终于要思想，这时日毫无益处，但在将来，你无需生活的技巧就有希望。诚然，那些被夺去的小孩子渴望活着。再者，那些夭折的年轻人，或许正打算老去，他们自己准备享受欢乐的日子；然而我们却不愿放下世上的一切。我曾以乖谬的心拖延，认为这世界的生命是真实的；我判断死亡会像你所想的一样到来——即一旦生命离去，灵魂也就死了并消失。但这些并非如此；天地的创始者与造主确已为兄弟所杀的兄弟要求追讨。不虔敬的人，主说：\BibleQuote{你的兄弟在哪里？}他否认了。因为\BibleQuote{你的兄弟的血从地上向我喊冤}。我看见你受苦，当你以为感觉不到什么时；但他活着，并占据了右边的位置。他享受着你这恶人所失去的喜乐；当你叫回世界时，他也已先行而去，并将不朽：因为你将在地狱哀号。的确，天主活着，他使死人复生，为给无辜和善良的人应得的赏报；但对凶暴和不虔敬的人，则是残酷的地狱。开始吧，你这被迷惑的人，要认识天主的审判。\par

% structure: p0057 source=h2 target=section reason=relative-heading-rank; base=h2

\section{第二十七章 愚拙的人哪，你并非对天主而死。}

愚拙的人哪，你并非完全死去；死时也不能逃脱那至高者。纵使你安排死后毫无知觉，你仍将被愚拙地胜过。创造世界的天主活着，他的律法宣告死者是存在的。但当你鲁莽地寻求没有天主而活着时，你判断死亡是消灭，并以为这是绝对的。天主并未如你所想那样安排，使死者忘记先前所行的事。如今治理者为我们造了死亡的居所，我们将在灰烬之后看见它们。你这愚拙的人被剥去了，你自以为死亡之后你就不存在，并使你的统治者和主宰无能为力。但死亡并非纯粹的虚空，如果你在心中再思量。你可以知道他是可渴望的，因为你将在他以后才认识他。你是肉体的统治者；诚然，肉体并未统治你。脱离肉体后，前者被埋葬；你在这里。凡人从肉体分离是正当的。因此，凡人的眼睛无法与（神圣之事）等同。我们的深奥如此使我们不能知晓天主的奥秘。现在，当你在软弱中死去时，应把尊荣归给天主，并相信基督将把你从死者中带回来。你应当在教会中赞美全能者。\par

% structure: p0059 source=h2 target=section reason=relative-heading-rank; base=h2

\section{第二十八章 义人复活。}

正义与良善、和平与真实的忍耐、以及对自己行为的关心，使人死后得生。但诡诈的心、恶意的、不忠的、邪恶的，逐渐毁灭自己，并滞留在残酷的死亡中。恶人哪，现在听听你因恶行得到什么。看看地上的审判官，他们现在在身体上以可怕的刑罚折磨人；要么为应得的人用刀剑预备惩罚，要么在长期的监禁中哀哭。你这最后的人，竟希望嘲笑那位创造万有的天主和掌管诸天的统治者吗？你发怒，你疯狂，现在你除去天主的名字，然而你仍不能逃脱他；他将按你的行为施行惩罚。现在，我愿你谨慎，免得你落入火焚之中。立即把自己交托给基督，好使良善伴随你。\par

% structure: p0061 source=h2 target=section reason=relative-heading-rank; base=h2

\section{第二十九章 致邪恶不信的富人。}

富人哪，由于贪婪地过分注视你所有的财富，你将挥霍那些你仍想紧抓不放的东西。你说：我不希望死后在这样的东西中活着。对伟大天主忘恩负义的人哪，你竟自视像神一样；他却在你一无所知时创造了你，后来又养育了你。他管理你的草地；他管理你的葡萄园；他管理你的牛群；他管理你所拥有的一切。你也不留意这些；或许你统治一切。那创造天空、土地和咸海的主，定意在黄金时代将我们本身归还给我们。只有你相信，你才能活在天主的奥秘中。认识天主，愚拙的人，他愿意你不朽，好使你在斗争中向他献上永恒的感谢。他自己的法律教导你；但因你寻求游荡，你不信一切，从此你要进入地狱。不久你放弃生命；你将被带到那个使你悲伤的地方：在那里，永恒的属灵惩罚将被承受；那里常有哀号；你在其中并非完全死去——在那里，你最终将太迟地宣告全能的天主。\par

% structure: p0063 source=h2 target=section reason=relative-heading-rank; base=h2

\section{第三十章 富人们，要谦卑。}

学习吧，你们这些即将死去的人，要向众人显示你的善。为何在众人中间，你要使自己显得与\Emph{你所是}不同？你往你不知道的地方去，又从那里无知地离去。你邪恶地对待你的身体；你永远渴求财富。你把自己抬得太高；你心怀骄傲，不愿看顾穷人。如今你甚至不供养在你权下的父母。唉，可怜的人，愿普通人远远逃离你们。他活着，我却毁灭了他；穷人大喊\OriginalTerm{我找到了}{\GreekText{εὕρηκα}}。不久你将被卡律布狄斯的愤怒驱赶，当你自己灭亡时。这样，你们富人是无纪律的，你们给别人立法，自己却不预备。脱下吧，远离天主的富人，脱去这样的邪恶，如果确实，或许，你所见到的能帮助你。趁你还有时间，做天主的侍从。正如榆树爱葡萄树，所以你们要爱卑微的人。现在，无结果实的人，要遵守对恶人可怖、对善人同样仁慈的法律；在顺境中要谦卑。富人们，除去欺诈的心，拿起平安的心。并要看你的恶行。你行善吗？我在这里。\par

% structure: p0065 source=h2 target=section reason=relative-heading-rank; base=h2

\section{第三十一篇。致审判官。}

你们所有审判官，要思想撒罗满的言语；他如何一句话就贬低了你们。礼物和赠品如何败坏审判官，由此，由此产生法律。你们永远喜爱送礼的人；当有案件时，不义的理由便得胜。这样我是无辜的；我一个卑微的人并不控告你们，因为撒罗满公开提出这亵渎。但你们的神即是你们的肚腹，赏报即是你们的法律。宗徒保禄提示这一点，我并不诡诈。\par

% structure: p0067 source=h2 target=section reason=relative-heading-rank; base=h2

\section{第三十二篇。致自娱者。}

如果地方或时间有利，或者人已晋升，那么便有一位新审判官。为何你现在因此而自高？未受教者，你亵渎那赐你生命之恩的主。在如此软弱中，你从未顾念祂。在晋升和收益中，你贪婪地依赖命运。对你没有法律，你也不在顺境中辨别自己。即使他们被算为金子，让风笛的曲调总是狂乱。如果你没有朝拜主的十字架，你便灭亡了。现在地方、时机和人都给了你，如果你相信；但若不，你将在祂面前恐惧。使你自己服从基督，把你的颈项放在祂之下。尊荣和一切事物的依靠都属于祂。当时机奉承你时，要更加谨慎。你未能如应当的那样预见命运的最后裁决，你将永远不能没有基督而再活。\par

% structure: p0069 source=h2 target=section reason=relative-heading-rank; base=h2

\section{第三十三篇。致外邦人。}

啊，凶猛无牧的人民，如今至少不要再漂泊。因为我这劝诫你们的人，也曾同样无知、漂泊。现在，因此，要取你们主的样式。举起你们狂野而粗硬的心。坚定地进入你们林中牧人的羊圈，在君王的屋顶下免受强盗之害。林中有狼；因此要躲进洞穴。你们作战，你们疯狂；你们看不见你们所在之处。要相信独一的天主，使你们死后能活，并在祂的王国中复活，那时将有义人的复活。\par

% structure: p0071 source=h2 target=section reason=relative-heading-rank; base=h2

\section{第三十四篇。此外，致无知的外邦人。}

未驯服的颈项拒绝负劳苦的轭。然后它喜欢在富饶的平原以草果腹。然而有用的母马不情愿地被制伏，当它初次被带入顺服时变得更不凶猛。啊人民，啊人，你的兄弟，不要成为野蛮的畜群。最终你自拔出来，你自己退隐。你确实不是家畜，不是野兽，而是生为人。你要明智地制伏自己，进入武装之下。你们跟随偶像的，不过是时代的虚妄。你们琐碎的心在几乎得自由时毁灭你们。那里有金子、衣服、银子被带到肘边；那里有战争；那里有情歌代替圣咏。当你们玩耍或期盼这样的事时，你们以为是生命吗？无知的人，你们选择已灭绝之物；你们寻求金器。这样你们不能逃脱瘟疫，虽然你们自己是属神的。你们不寻求天主差遣到地上来让人阅读的恩宠，反而如野兽般漂泊。如果你们相信，那先前说的黄金时代将临到你们，你们将再次开始永远活出不死生命。那也被允许知道你们从前是什么。你们要顺服于统管万有的天主。\par

% structure: p0073 source=h2 target=section reason=relative-heading-rank; base=h2

\section{第三十五篇。论生命树与死亡树。}

亚当是第一个跌倒的，并且为逃避天主的诫命，贝里雅耳借着棕榈树的欲望诱惑了他。他所行的，无论善恶，也传给了我们，作为从他而生的一切人的首位；因此我们因他而死亡，正如他自己，远离神圣，成了圣言的弃儿。当六千年完成时，我们将是不死的。苹果树被尝，死亡进入了世界。借着这死亡树，我们为来生而生。那结出果实的生命依赖于树——即诫命。现在，因此，要怀着信心摘取生命的果实。一条法律从树上传下，为原始人所畏惧，由此因忽略起初的法律而死亡。现在伸出你的手，取生命树的果实。随后而来的主的卓越法律，也出自这树。最初的法律丧失了；人从能取的取食，他朝拜被禁的神明，生命的邪恶喜乐。拒绝这种分享；知道它应当是什么就够了。如果你愿意活着，就顺服于第二法律。要避免神庙的敬拜、魔鬼的神谕；转向基督，你们将成为天主的同伴。天主的法律是神圣的，它教导死人活过来。唯独天主命令我们向祂献上赞美的诗歌。你们所有人都要彻底避开魔鬼的法律。\par

% structure: p0075 source=h2 target=section reason=relative-heading-rank; base=h2

\section{第三十六篇。论十字架的愚拙。}

我已论及死亡源自的双重标记，又说生命常由此而出；但那十字架，对奸淫的百姓成了愚妄。可怕的永恒之王藉十字架预示\Emph{这些事}，好使他们如今相信祂。愚昧人哪，活着竟在死亡中！加音以邪恶的发明杀了他的弟弟。由此，恩诺士的儿子们被称为加音的后代。于是那从不将灵魂归于天主的邪恶百姓在世上增多。相信十字架成了可怖之事，他们却说自己是正义地生活。第一条律法在于树；第二条也是如此。而第二条律法首先以平安胜过了那可怖的律法。被高举起来，他们却冲入虚妄的曲解中。他们不愿承认那被钉子刺透的主；但当祂的审判来临，那时他们将辨识祂。但亚伯尔的后裔已经相信了那仁慈的基督。\par

% structure: p0077 source=h2 target=section reason=relative-heading-rank; base=h2

\section{XXXVII.　犹太化的狂热者。}

怎么？你是半个犹太人？你愿意半是亵渎吗？由此你死后将不能逃脱基督的审判。你自己盲目地游荡，又愚蠢地混入盲人之中。于是盲人领盲人，双双掉入坑里。你往你不晓得的地方去，又从那里无知地退离。让学习的人去就学者，让学者离去。但你却往那些你不能学到什么的人那里去。你行至门前，又从那里走向偶像。首先问问律法中命令了什么。让他们告诉你，是否命令敬拜诸神；因为他们就自己特别能做的事而言，是被忽略的。但因他们犯下那同一罪行，他们不讲论天主的诫命，只讲论奇妙的事。然后，他们却盲目地领你一同跌入坑里。有死亡为他们所熟知而无需讲述，或因犁的堆叠封闭了田地。全能者不愿他们认识他们的君王。为何如此邪恶？祂自己逃离了那些血腥的人。祂藉增补的律法将自己赐给我们。从此他们如今与我们一同隐藏，被他们的君王遗弃。但你若以为在他们那里有希望，你全然错了，如果你敬拜天主又敬拜异教庙宇。\par

% structure: p0079 source=h2 target=section reason=relative-heading-rank; base=h2

\section{XXXVIII.　致犹太人。}

常是邪恶、悖逆、硬着颈项的，你不愿被克服；如此你将成为继承者。依撒意亚说你们是心硬的。你们看着梅瑟在愤怒中摔碎的律法；而同一主赐给了他第二条律法。他在那里寄望；但你们，半得医治，弃绝了它，因此你们不配进入天国。\par

% structure: p0081 source=h2 target=section reason=relative-heading-rank; base=h2

\section{XXXIX.　同样致犹太人。}

请看肋阿，她是会堂的预像，雅各伯以眼目如此昏花中接受了她为标记；然而他又为那被爱的年幼者服事：一个真实的奥秘，我们教会的预像。试想从天上论及黎贝加所说的大量话语；由此，效法那外方人，你们可以信从基督。然后来到塔玛尔和双胞胎的子孙。看加音，第一个耕地的，和亚伯尔那牧羊的，他是在兄弟的毁灭中无玷的献祭者，被他的兄弟所杀。如此，你们当明白：年幼的蒙基督悦纳。\par

% structure: p0083 source=h2 target=section reason=relative-heading-rank; base=h2

\section{XL.　再次致同样的人。}

没有像你们这样不信的百姓。邪恶的人哪！在如此多的地方，并如此多次被那大声呼喊的律法所责斥。崇高者藐视你们的安息日，并完全弃绝你们按律法的一切每月节期，使你们不可向祂献命定的祭；那告诉你们投石以赎过犯的。若有人不信祂死于不义之死，而所爱者因别的律法得救，从此生命被悬于树上，却仍不信祂。天主自己就是生命；祂自己为我们被悬。但你们以硬心侮辱祂。\par

% structure: p0085 source=h2 target=section reason=relative-heading-rank; base=h2

\section{XLI.　论假基督的时代。}

依撒意亚说：这就是那震动大地和如此多君王的人，在他之下土地变为荒芜。听先知如何论及他预言。我并未详述，而是粗略提及。那时，无疑地，当他出现，世界将终结。他自身将把大地分为三股统治势力，此时，尼禄将从地狱被举起，厄里亚将先来封印那些蒙爱者；对此，非洲地区和北方民族，整个大地四面八方，将战栗七年。但厄里亚将占据一半时间，尼禄占据另一半。然后那巴比伦淫妇化为灰烬，其余烬将从那里向耶路撒冷推进；那拉丁征服者于是说：我是基督，你们常向我祈祷的；而那起初受骗的便联合起来赞美他。他行许多奇事，因为他的假先知。尤其为使众人相信他，他的偶像会说话。全能者赐予它能力显得如此。犹太人从他那里重述圣经，同时向至高者呼喊说他们受了骗。\par

% structure: p0087 source=h2 target=section reason=relative-heading-rank; base=h2

\section{XLII.　论全能基督、永生天主那隐藏而圣洁的百姓。}

愿那隐藏的、终末的、圣洁的子民被渴望；并且，诚然，愿我们不知道它居于何处，因那九个半支派行事……他曾吩咐他们按从前的法律生活。如今，我们全都生活：法律的传授是新的，正如法律本身所教导的，我更清楚地指示你们。两个半支派被留下：为何那半支派要从他们中被\Emph{分离}出来？为使他们成为殉道者，当祂向祂在世上蒙选的人发动战争时；或者，毫无疑问，圣先知们的歌队将一同兴起攻击那百姓，他们本该遏制那些被猥亵的马匹用踢踏的后蹄所残杀的人；并且那队列也决不会在任何时候轻率地冲向\Emph{和平的恩赐}。那些支派被撤回，而基督的一切奥秘由他们在整个时代中完成。再者，他们从两兄弟的罪行中兴起，在他们的征兆下追随了罪行。这些血腥者如此被分散并非不当：他们将再次为基督的奥秘而聚集。但随后，法律中所预表的事正急速趋向完成。全能基督降临到祂的选民那里，他们长久以来被遮蔽在我们的视线之外——他们已成为如此成千上万——那便是真正的天主子民。那时，儿子不会先于父亲而死；他们身体也无疼痛，鼻腔也无息肉。那些离世的人在年迈时安卧于床榻，成就了法律的一切，因此他们得到保护。他们被命令从主的右边经过；当他们如前经过时，祂使河流干涸。主自己也照样与他们一同前行。祂已来到我们这一边，他们与天国的君王同来；在他们的旅程中，我该述说天主将成就何等事呢？群山在他们面前下沉，泉水涌出。受造物欢喜看见天国的子民。然而在此，他们急速去保护那被掳的妇人。但那占有她的邪恶君王，一听见便逃往北方，并召集所有\Emph{他的随从}。再者，当那暴君自冲于天主的军队时，他的兵士被天上的恐怖击溃；那假先知自己也与那恶者一同被主谕旨擒获；他们被活生生地交付于\OriginalTerm{革哈那}{Gehenna}。首领和领袖被命令从他那里服从；然后圣洁者将进入他们古老母亲的怀抱，使那些被他邪恶说服的人也得以安息。他以各种刑罚折磨那些信赖他的人；他们来到结局，罪过便从世上被除去。主将开始以火施行审判。\par

% structure: p0089 source=h2 target=section reason=relative-heading-rank; base=h2

\section{XLIII. 论今世的终结}

号角在天上发出信号，狮子被带走，忽然间天宇发出轰鸣而黑暗降临。主垂下眼目，大地震动。祂呼喊，使普世都能听见：\Quote{看，我长久静默，在这期间容忍了你们的行径。}他们一同呼喊，抱怨并呻吟，却为时已晚。他们哀号，他们痛哭；恶人找不到立足之地。母亲能怎样对待吃奶的婴儿，当她自身被焚烧时？在火焰中，主将审判恶人。但火不会触及义人，却必全然舔舐他们。他们在一处滞留，但有一部分在审判中哭泣了。如此炎热，连石头都要熔化。风聚集成闪电，天上的义怒肆虐；恶人无论逃往何处，都被这火抓住。没有救助，也没有海船。\OriginalTerm{阿们}{Amen}火焰临到列国，玛待人和帕提亚人焚烧一千年，正如若望隐藏的话所宣告的。因为此后一千年，他们被交付于\OriginalTerm{革哈那}{Gehenna}；而他们的作为者，也同他们一起被焚烧。\par

% structure: p0091 source=h2 target=section reason=relative-heading-rank; base=h2

\section{XLIV. 论第一次复活}

那城将在第一次复活中从天而降；这便是我们所能述说的如此天工的建造。我们要向祂复活，我们这些曾奉献于祂的人。他们将成为不朽坏的，甚至已活在没有死亡中。在那城里不再有任何忧伤，任何呻吟。那些在\OriginalTerm{假基督}{Antichrist}手下战胜了残酷殉道的人也要来，他们自己活整个时期，领受祝福，因为他们曾遭受恶事；并且他们自己嫁娶，生育一千年。地上的一切出产都已预备好，因为大地被更新，源源不绝地丰产。那里没有雨水；没有寒冷进入黄金的营寨。不像如今，没有围攻，没有掠夺，那城也不渴求灯烛之光。它从自己的建造者发光。再者，它服从祂；宽一万二千斯达狄，长与高亦然。它的根基奠于地上，却将头高举至天。在城门前的城中，日与月照耀；恶人被围困在痛苦中，为义人的滋养之故。但一千年后，天主将消灭所有那些邪恶。\par

% structure: p0093 source=h2 target=section reason=relative-heading-rank; base=h2

\section{XLV. 论审判日}

因不信者之故，我再添加一些关于审判日的话。再次，主所发出的火将被指定。大地发出真实的呻吟；那时那些在最终行程中的人，以及所有不信者，都\Emph{呻吟}。整个自然被火焰转变，那火却避开祂圣徒的营寨。大地从根基被焚烧，山岳熔化。海洋无物存留：它被强大的火征服。此天消逝，星辰和这些事物都被改变。另一种新的天空和永存的大地被铺设。然后，那些配得的人被送离，处于第二次死亡中；但义人被安置在内室中。\par

% structure: p0095 source=h2 target=section reason=relative-heading-rank; base=h2

\section{XLVI. 致候洗者}

简言之，我劝告所有为你们的救恩而离弃偶像的、在基督内的信徒。起初，若你们以任何方式陷入谬误，仍当应召为基督舍弃一切；既已认识天主，就当成为良善而可嘉许的新兵，并让童贞的谦逊与你们一同持守纯洁。心智当警醒于善事。谨防重陷昔日的罪恶。在圣洗中，你们诞生的粗劣衣裳被洗净。若有任何有罪的慕道者受罚记，就让他生活在基督信仰的标记中，虽非毫无损失。你们全部的要点即在于：你们务必永远远离严重的罪。\par

% structure: p0097 source=h2 target=section reason=relative-heading-rank; base=h2

\section{XLVII. 致信友}

我劝告信友不可怀恨你们的弟兄。仇恨被殉道者视为不虔敬，竟至于火焰。殉道者若其宣认属于此类，便遭毁灭；亦未教导说邪恶可由流血而赎。法律赐予不义之人，使他自制。因此他当远离诡诈；你们也当如此。若你们将争端延伸至弟兄，便是双重得罪天主；由此你们无法避免随从先前行径的罪恶。你们既已领受圣洗一次，还能再次浸入吗？\par

% structure: p0099 source=h2 target=section reason=relative-heading-rank; base=h2

\section{XLVIII. 信友啊，谨防邪恶}

飞鸟受骗，林中的野兽在林中受那些始终导致其毁灭的诱惑所欺，洞穴与食物在它们追随时欺骗它们；它们不知如何躲避邪恶，也不受法律约束。法律赐予人，并赐予应择取的生命教义，人由此牢记能谨慎生活，回想自己的位置，并除去属于死亡的事物。凡离弃规范者，严厉谴责自己：或为铁链所缚，或从自己的品级被贬；或丧失性命，失去应享之物。被榜样警戒，不可严重犯罪；经洗礼池更新，倒要拥有爱德；远远逃离捕鼠夹的诱饵，那里有死亡。许多殉道无须流血。不贪求他人之物；渴望获得殉道的益处；管束舌头，你们当自谦；不乐于使用暴力，也不以暴力还报所施于你们的暴力，你们将有一颗忍耐的心，明白你们是殉道者。\par

% structure: p0101 source=h2 target=section reason=relative-heading-rank; base=h2

\section{XLIX. 致补赎者}

你已成为补赎者：日夜祈祷；然而不可远离你的母亲——教会，至高者方能怜悯你。你罪过的告解必不落空。同样，在你受控告的状况中，学习公开哭泣。你若有了伤口，就寻求草药和医生；并且你仍能在惩罚中减轻你的痛苦。因我甚至要承认，在此我独自一人，而恐惧必须被抛在身后。我亲身感受过毁灭；因此我劝告受伤者更谨慎行走，将你们的头发和胡须置于地上的尘土，身穿麻衣，向至高君王恳求，祂必帮助你们，免得你们或许从百姓中灭亡。\par

% structure: p0103 source=h2 target=section reason=relative-heading-rank; base=h2

\section{L. 凡背弃天主者}

再者，当战争进行，或仇敌来袭时，若一人能或征服或隐藏，便是巨大的胜利；但被他们所擒的人将是不幸的。他失去国家和君王，因他不愿为真理、为国家或为生命而英勇作战。他宁死也不愿受制于蛮族的君王；而甘愿将自己转移到无法无天的仇敌那里的人，当寻求奴役。那时，若你在作战中为君王而死，你已得胜；或若你交出双手，你已未受法律损害而灭亡。仇敌渡河；藏身于你的隐匿处；或不论他能否进入，都不要逗留。处处确保自身和朋友的安全：你已得胜。并要警惕，免得任何人进入那隐匿处。若有人向仇敌暴露自己，便是可耻的事。那不知如何征服、跑去投降的人，既非为自身亦非为国家的利益，怯懦地放弃了赞美。那时他本不愿活着，因为生命本身也将逝去。若有人没有天主，或来自仇敌的亵渎者，他们便成了鸣的锣，或聋如蝮蛇：这样的人理应多多祈祷或隐藏自己。\par

% structure: p0105 source=h2 target=section reason=relative-heading-rank; base=h2

\section{LI. 论婴孩}

仇敌突然以战争淹没我们；在他们能逃跑之前，他已掳掠了无助的孩童。他们虽被视为被俘，却不可受责备；我实在也不为他们开脱。或许他们因父母的过犯而该当如此；因此天主舍弃了他们。然而，我劝告成年者拿起武器，并好像从母腹重生到母亲（教会）那里。让他们避开那可怕、永远血腥、不虔敬、难驯服、活得像野兽的法律；因为当另一场战争偶然要进行时，那应能征服甚至应能正确知道如何警惕的人。\par

% structure: p0107 source=h2 target=section reason=relative-heading-rank; base=h2

\section{LII. 逃兵}

因为逃兵并非全属同一类。一人是邪恶的，另一人部分退却；然而对两者都判定了真正的审判。因此，基督被对抗，正如凯撒被服从。你若曾为罪犯，当寻求君王的庇护。向祂恳求；俯伏于祂，向祂告解：祂将赐予一切，因为万有也属于祂。营地已重建，当心别再犯罪；不要像士兵长久在野兽的洞穴中游荡。在你们停止无度行为上，成为罪恶吧。\par

% structure: p0109 source=h2 target=section reason=relative-heading-rank; base=h2

\section{LIII. 致基督的士兵}

当你报名参加这场战争时，你已被缰绳约束。因此开始摒弃你从前的行为。躲避奢华，因为劳苦是威胁的武器。你必须以全部德行服从君王的命令，若你渴望在喜乐中抵达末时的境界。他是一名好士兵，总在等待那将要享受的事物。不要自满，彻底摒弃懈怠，以便每日为摆在你面前的事做好准备。事先谨慎，在早晨再次检视军旗。当你看见战争时，争取最近的搏斗。这是君王的荣耀：看见军队准备就绪。君王就在眼前，愿你奋勇战斗超越他的期望。他预备了赏赐。他欣然期待胜利，并指派你为适宜的追随者。除了为\OriginalTerm{贝利亚尔}{Belial}，不要吝惜自己；反要勤勉，好让他为你的死亡赐予名声。\par

% structure: p0111 source=h2 target=section reason=relative-heading-rank; base=h2

\section{五四 论逃亡者}

那些丧亡者的灵魂理当自行分离。由他所生，却又回归属于他的事物。\OriginalTerm{加音}{Cain}的根，受诅咒的苗裔，迸发出来，在野蛮的君王之下逃入为奴的民族；在那定夺之日，永火将折磨他们。逃亡者将漫无纪律地游荡，不受法律约束，游走于道路的隘口。因此，这些人正是刑罚未能遏制的那种。若他们不愿存活，理应为偶像所看见。\par

% structure: p0113 source=h2 target=section reason=relative-heading-rank; base=h2

\section{五五 论莠子的种子}

论那混在教会中的莠子的种子。当收割的时期满足，长出的莠子便从果实中分离出来，因为天主并未派遣它们。农夫收集所有那些莠子，将它们分离。法律是我们的田地；凡在其中行善的，统治者本人必赐予真正的安息，因为莠子要用火焚烧。因此，若你以为它们在同一个之下被延迟，你就错了。我称你们为不结果实的基督徒；那无花果树的诅咒，因主的话语而无果，立刻枯萎了。你们不做善工，不为宝库预备礼物，却又如此徒然以为能得主的欢心。\par

% structure: p0115 source=h2 target=section reason=relative-heading-rank; base=h2

\section{五六 致伪善者}

你竟对那如此公开宣告、在如此众多先知的天上话语中呼喊的法律装作不知吗？若一位先知只是向着云彩呼喊，由他发出的主的话语必定已经足够。主的法律自身宣告于如此众多的先知书卷中；无一人为邪恶开脱；因此你甚至从心里渴望看见好事，也寻求以欺骗度日。那么，为何法律本身要如此费力地颁布呢？你滥用主的诫命，却仍称自己为他的儿子。你已被看见，若你毫无理由地如此行事。我说，全能者寻求谦逊的人作他的儿子，那些心地正直、忠于神圣法律的人；但你已经知道他已将恶人投于何处。\par

% structure: p0117 source=h2 target=section reason=relative-heading-rank; base=h2

\section{五七 论世俗之事必须完全避免}

若某些教师，因贪图你们的礼物或畏惧你们的人格，而对你们每件事放宽要求，我不仅不悲伤，还被迫说出真理。你们随同恶者的群众前往虚浮的表演，那里\OriginalTerm{撒殚}{Satan}在竞技场中喧嚣作工。你自以为凡令你喜悦的一切都是合法的。你是至高者的苗裔，却与魔鬼的子女混杂。你愿看见你先前所弃绝的事物吗？你再次与它们交往？那受傅者于你何益？或者，因软弱，你愚昧地亵渎，这被允许吗？不要爱世界，也不要爱其中的事物。这是天主的圣言，而你看为美善。你遵守人的命令，却躲避天主的。你信赖那礼物，教师因此闭口不言，不再告诉你神圣的诫命；我既然说出真理，你理当仰望至高者。你当将自己交托为那位你曾是儿子者的追随者。若你寻求生活，作为一名有信仰的人，却像外邦人一样行事，世界的欢乐将你从基督的恩宠中移开。你以无纪律的心思寻求那些你自以为轻易合法的事，既是你亲爱的伶人及其乐曲；你也不关心这般人的后代应当嘀咕愚行。你以为自己在享受生命，却粗心地犯错。至高者命令，而你躲避他正义的诫命。\par

% structure: p0119 source=h2 target=section reason=relative-heading-rank; base=h2

\section{五八 论基督徒应该如何}

当主说人该叹息着吃饼，你在此正在做什么？你渴望喜乐地生活！你寻求废除至高天主在最初形成人时宣判的判决；你希望抛弃法律的约束。若全能天主命令你以汗度日，你活在享乐中，便已与他疏远。圣经记载主曾向犹太人发怒。他们的儿子，因食物得以焕发，起来玩耍。现在，我们为何追随这些受割礼的人呢？他们在哪方面灭亡，我们应当警惕；你们中的大部分，沉溺于奢华，听从他们。你以染料玷污自己，而违犯法律；宗徒向你呼喊，是的，天主藉着他呼喊。他说：你的放荡自身毁灭你。因此，要成为基督所愿你成为的那样：温良，并在他内喜乐，因为在世界上你是忧愁的。奔跑，劳苦，流汗，忧愁作战。希望伴随劳苦而来，棕榈枝赐予胜利。若你渴望得振奋，就帮助并鼓励殉道者。等待那在死亡通道中来临的安息。\par

% structure: p0121 source=h2 target=section reason=relative-heading-rank; base=h2

\section{五九 致永生天主的教会的已婚妇女们}

你，基督徒妇女，希望主妇们如同世间的贵妇一般。你用黄金或朴素的丝绸衣服装扮自己。你将法律的恐吓从耳中抛向风中。你用魔鬼的一切奢华炫耀虚荣。你在镜子前梳妆，卷曲的头发从额前向后翻。此外，你怀着恶意，在你纯净的眼睛上涂抹假药、\OriginalTerm{锑粉}{stibium}，以描画的美貌，或将头发染黑。天主是俯察各人心的监督者。但这些事对端庄的妇女并无必要。要以贞洁和端庄的情感刺透你的心。天主的法律证明，这样的法律从相信的心中消失；对于一位被丈夫认可的妻子，让她满足于这一点，不是靠她的衣着，而是靠她的良好性情。穿上因寒冷、炎热或烈日所需的外衣，只为使你被认可为端庄，并在天主子民中展示你才能的恩赐。你从前极为显赫，如今要自甘卑微。那位躺卧无生命的妇人，因寡妇们的祈祷而复活。她理应如此，从死亡中复活，不是靠她昂贵的衣服，而是靠她的施舍。善心的主妇们，逃避虚荣的装饰；这样装扮适合那些流连妓院的女人。战胜那恶者，基督的端庄妇女们。在施舍中显出你的一切财富。\par

% structure: p0123 source=h2 target=section reason=relative-heading-rank; base=h2

\section{六十、致同一人}

听我的声音，你这位希望保持基督徒妇女身份的人，听那有福的\Person{保禄}如何吩咐你装扮。此外，从天上发言的导师和作者\Person{依撒意亚}，因他厌恶那些追随世俗邪恶的人，说：\Quote{熙雍的女子高傲的，必被降卑。}在天主眼中，信主的基督徒妇女不应装饰是不对的。你，奉献于天主的人，要效仿外邦人的方式出去吗？天主的使者，在法律中大声呼喊，谴责这样的不义妇女，她们如此装饰自己。你染你的头发；你用黑色描绘你的眼睑；你一缕一缕地把你漂亮的头发放在涂了粉的额头上；你用某种红色涂抹你的脸颊；此外，耳环沉重地垂着。你用项链埋没你的脖子；用宝石和黄金捆缚配得上天主的手，这是不祥的预兆。我何必提你的衣裙，或魔鬼的全部奢华？当你想要取悦世界时，你就在拒绝法律。你在家中跳舞；代替圣咏，你唱情歌。你，虽然可能贞洁，却因追随恶事而证明自己并非如此。因此基督使你这样的人与外邦人同等。要取悦那被歌颂的唱诗班，并以炽热的爱向已平息怒火的基督欣然献上你的馨香。\par

% structure: p0125 source=h2 target=section reason=relative-heading-rank; base=h2

\section{六十一、在教会中向全体天主子民讲话}

弟兄们，我并非义人，从污秽中被举起，我也不自高；但我为你们忧伤，因为看到如此众多的人中，没有一人在竞赛中得冠冕；当然，即使他自己不争战，至少也应鼓励别人。你责备灾难；哦肚腹，你以奢侈填满自己。弟兄在对抗世界的战斗中有武器；而你，满有财富，既不争战，也不在他争战时站在他身边？愚昧的人，你难道不知一人是为许多人争战吗？若他得胜，整个教会都系于他。你看见你的弟兄被阻碍，他与敌人争战。你在营中渴望和平，他在外面拒绝和平。要怜悯，使你凡事得救。你也不畏惧主，他以如此宣告呼喊；祂甚至吩咐我们也要给敌人食物。期望从那位\Person{多俾亚}而来的你的饭食，他每日总是把全部饭食与穷人分享。你寻求喂养他，愚昧的人，他却再喂养你。你希望他为我备办，而他却在为自己预备葬礼？被贫乏压迫、几乎憔悴的弟兄向那丰食饱腹、肚腹鼓胀的人呼喊。关于主日你怎么说？若他尚未预先到场，从人群中召出一个穷人，带他赴你的宴席。在经书中，因基督得到慰藉，你的盼望所在。\par

% structure: p0127 source=h2 target=section reason=relative-heading-rank; base=h2

\section{六十二、致渴望殉道者}

儿子，既然你渴望殉道，请听。要像\Person{亚伯尔}，或像\Person{依撒格}本人，或像\Person{斯德望}，他为自己选择了义人的道路。你确实渴望那是属于有福之人的事。首先，以你的善行战胜那恶者，过正直的生活；当你的君王看见你时，你就安全了。这是祂自己的时候，我们为两者活着；这样，即使战争结束，殉道者也将平安地进入。许多人确实犯错，他们说：\Quote{我们用血战胜了那恶者；}若他仍存留，他们便不愿得胜。他埋伏着灭亡，恶人如此感受；但合法的人并不感受施加的刑罚。你要呼喊，急切地用拳头捶打你的胸膛。即使现在，你若以善行战胜，你在祂内就是殉道者。因此，你用言语渴望高举殉道者，在平安中以善行披上自己，并要安稳。\par

% structure: p0129 source=h2 target=section reason=relative-heading-rank; base=h2

\section{六十三、每日的战争}

你想发动战争，愚昧的人，仿佛战争在和平中。从起初受造的日子起，你最终争战。情欲催迫你，就有战争；要与它作战。奢侈引诱你，要拒绝它；你就战胜了战争。要节省酒的丰足，免得藉酒误事。制止你的舌头不说咒骂的话，因为你用它朝拜主。抑制怒气。对众人和平。当心不要践踏比你卑微、被苦难压伤的人。只把自己当作保护者，不伤害人。引导自己走在正直的道路上，不被嫉妒玷污。在你的财富中，对卑微的人要温和。付出你的劳动，给赤身者衣服。如此你将得胜。不为任何人设陷阱，因为你服事天主。看起初，嫉妒的仇敌从那里灭亡了。我不是教师，但法律本身藉其宣告教导。你空说如此大的话，却想在一瞬间不经劳动就为基督树立殉道的事迹。\par

% structure: p0131 source=h2 target=section reason=relative-heading-rank; base=h2

\section{六十四、论贪欲的热忱}

存有贪欲，你便由此丧亡，因你燃烧着对邻人的妒火。你内在燃起忿怒时，便自我毁灭。嫉妒的人啊，你妒忌那与邪恶抗争的他人，渴望同样获得那么多财富。当你试图攻击他时，法律并不如此看待他。你依赖万事万物，活在贪婪中；虽然你自知有罪，却以自己的判断定己罪。眼目的贪求永不止息。所以，你若回头思量，便知贪欲是虚妄……因此天主呼喊说：\Quote{愚昧人，今夜就要索回你的灵魂。}死亡紧随你。那些塔冷通将归于何人？将不义之财藏入隐秘的库中，而上主赐予每人每日的生计。让别人积累吧，你当寻求善生。当你的良心在天主面前无愧，你将战胜一切；但我并非说你要在众人前夸耀，当你无欺诈度日、等候你的日子时。鸟儿在食物中丧亡，或疏忽中粘在捕鸟胶上。要以为你的单纯中有许多需要提防。让别人越界吧。你总是前瞻。\par

% structure: p0133 source=h2 target=section reason=relative-heading-rank; base=h2

\section{第六十五章  出自邪恶的施舍}

为何你愚昧地借他人的创伤伪充善人？你施舍之处，另有他人日日哭泣。你难道不信天主从天上看见这些？至高者说：\Quote{祂不悦纳恶人的礼物。}当你获得高位时，你便要欺压穷苦人。有人送礼是为了贬低他人；或者，你若放高利贷，取利百分之二十四，你便想行慈善以除净自己，以邪恶之物行邪恶。全能者绝对弃绝这样的行为。你施舍的\Emph{是}从泪水中强夺来的；那因不堪重利而穷困的候选人为此哀叹。此外，你为催债者制造了机会，你的仇敌如今便是百姓；你这献身者已因酬报而变成恶人。你还想以工资之利赎罪。邪恶的人啊，你欺骗自己，却骗不了别人。\par

% structure: p0135 source=h2 target=section reason=relative-heading-rank; base=h2

\section{第六十六章  论虚假的和平}

指定的时刻临近我们的人民；世上有和平，而同时，因世界的诱惑，毁灭正压垮我们（毁灭）那被你分裂为异派的鲁莽之人。要么服从城邑的法律，要么离开它。你看见我们眼中的木屑，却看不见自己眼中的梁木。虚假的和平正向你而来；迫害猖獗；伤口不显；如此，你未经杀戮而丧亡。战争在暗中进行，因为在和平之中，你们几乎无人谨慎行事。啊，可怜的、预定了宰割的人，你们赞美虚假的和平，对你们有害的和平。你们成了基督以外的士兵，已丧亡了。\par

我劝某些读者只要思量，并以生活的榜样给他人提供素材，避免纷争，躲避诸多争吵；压制恐惧，永不骄傲；此外，谴责恶人的伪义服从。使自己肖似你的基督师尊，小孩子们。因你们的善行而成为田野中的百合花；你们承载法令时便得福；你们是会众中的花朵；你们是基督的明灯。持守你们所是的，你们便能述说它。\par

% structure: p0138 source=h2 target=section reason=relative-heading-rank; base=h2

\section{第六十八章  致圣职者}

执事啊，要纯洁地施行基督的奥迹；因此，圣职者啊，执行你们师尊的命令；不要扮演正义审判的角色；以一切加强你们的职分，为有学问的人，仰望上界，永远忠于至高的天主。将祭台神圣的职事忠实地献给天主，在神圣之事上预备好作榜样；你们自己应向牧者俯首，如此你们方能蒙基督悦纳。\par

% structure: p0140 source=h2 target=section reason=relative-heading-rank; base=h2

\section{第六十九章  致天主的牧者}

牧者若曾宣认，便加倍了战斗。此外，宗徒吩咐这样的人应为教师。让他作有耐心的统治者；让他知道何时可以放松缰绳；让他先恐吓，再以蜂蜜涂抹；并让他首先自身遵行自己所说的。留恋世俗之事的牧人被看作有过失，无人敢当他的面说什么。\OriginalTerm{革赫纳}{Gehenna}本身在地狱中冒着谣言的气泡。那以犹疑的眉头摇摆的可怜人民有祸了！若有这样的牧人临到它，它几乎要灭亡。但虔诚的人约束它，治理得当。蜂群在合适的蜂王下喜乐；这样的人有希望，整个教会得存活。\par

% structure: p0142 source=h2 target=section reason=relative-heading-rank; base=h2

\section{第七十章  我向年长者发言}

时势要求我单独向你们讲真理。\par

他常因一句话受劝诫，而许多人拒绝。我希望你们单单向我来泄愤，使众人之心在诱惑者面前战栗。请思量那真正生恨的话语，（想想）我近来如何多次预言了一场欺骗的和平，而唉，诱人的恶者已悄然临到你们，因为你们不知其诡计迫近，你们已丧亡；你们行极其苦涩之事，但这正是世界的特征；任何你们为之说情的人，都不白白行动。那躲避你们火的人，跳入漩涡。然后那可怜的人被剥光，向你们求助。审判官们自己也为你们的欺诈战栗……若篇幅较短，我就不费如此多行了。你们这些教导的人，请留心你们情愿牧养的人，因为你们为自己接受宴席并享用它们。为了这些事，你们已几乎进入大地的根基。\par

% structure: p0145 source=h2 target=section reason=relative-heading-rank; base=h2

\section{第七十一章  探望病人}

如果你的弟兄软弱——我指的是穷人——不要空手去探望卧病的他。在天主内行善；以你的钱财尽服从。如此，他将康复；或若他去世，就让那无力回报你的穷人得到安抚，而那世界的创造主与作者会替他回报。你若不愿去看望穷人（他们总是令人厌恶），就送钱去，送些能使他自己康复的东西。同样，你若贫穷的姐妹卧病在床，让你们的妇人们开始给她送食物。天主亲自呼喊：\Quote{把你的饼分给有需要的人}。探望不在于言语，而在于施惠。你的弟兄因缺乏食物而生病，这是邪恶。不要只以言语满足他。他需要饮食。要立刻看顾那些的确衰弱、无法自助的人，立刻给他们帮助。我保证，天主将以四倍偿还给你。\par

% structure: p0147 source=h2 target=section reason=relative-heading-rank; base=h2

\section{第七十二章 对健康的穷人}

健康的贫穷若无财富能做什么？的确，你若有余力，就立刻也分施给你的弟兄。为一人对你自己负责，免得你被称为骄傲。我保证你将比富人活得更有平安。请侧耳倾听伟大撒罗满的教训：\Quote{天主憎恨穷人在高位上诉说。}因此你要谦卑，尊敬有权能者；因为温和的言语——你知道那句谚语——能融化人心。服侍能征服人，即使有积怨。若舌头沉默，你便一无所获。若没有合宜的技艺谋生，就当遵从那有权能者的命令，或给予援助，或加以指导。不要因健康的人有信德而感到羞耻或忧伤。此外，在库房里，你应当以自己的劳动奉献，正如那位基督所悦纳的寡妇一样。\par

% structure: p0149 source=h2 target=section reason=relative-heading-rank; base=h2

\section{第七十三章 不可哀哭儿子}

虽然儿子的死亡使人心存忧伤，但身穿黑衣或哀哭他们都是不应当的。主明智地说，你应当心中忧伤，而非外在表现，因为外在表现一周便完结。在撒罗满的书卷中，关于复活的主的应许已被你遗忘，你若使儿子成为殉道者，却又以口哀哭他们。你难道不羞于毫无节制地像外邦人那样哀哭儿子吗？你撕裂面容，捶打胸膛，脱去衣服；难道你不敬畏主，你渴望见到祂的国度？当按正理哀悼，但不要因他们而行恶。你就是这样的人。你比外邦人差多少？你如同出自魔鬼血统的群众。你哭喊说他们灭绝了。哦，虚伪的人，你灭亡有什么益处？父亲没有悲痛地将儿子带到祭台前杀死，先知也没有为死去的儿子悲痛哀哭，甚至没有父母哭泣。但那奉献于天主的人正在匆匆死去。\par

% structure: p0151 source=h2 target=section reason=relative-heading-rank; base=h2

\section{第七十四章 论殡葬排场}

你寻求在意死亡的排场，是在错误之中。作为天主的仆人，你即使在死亡中也应当中悦祂。哀哉，无生命的尸体在死后还要被装饰！哦，真正的虚妄，竟然渴望给死人荣誉！一颗被世界束缚的心灵，即使在死亡中也不献身于基督。你知道那些箴言。他想要被人抬着经过广场。同样，你这与他相似、以未经训练的心活着的人，希望在你死时有一个快乐蒙福的日子，好让人们聚集，使你看见赞扬伴随着哀悼。你没有预见你死后可能该去哪里。看，他们跟随你；而你或许已在燃烧，被驱向惩罚。排场对死人有什么益处？你会被指控，因为你因那些聚集而寻求它们。你渴望活在偶像之下。你欺骗自己。\par

% structure: p0153 source=h2 target=section reason=relative-heading-rank; base=h2

\section{第七十五章 致圣职者}

他们将在复活节——我们最蒙福的那日——聚集；让那些寻求神圣款待的人欢乐吧。让他们得到足够的供应，酒和食物。回顾源头，这些事如何能为你而说是出自于你。你缺乏一份献给基督的礼物，缺乏适度的花费。既然你们自己不做，又如何能说服这些人遵行律法的义，哪怕一年一次？因此许多人对你们起了亵渎的猜疑。\par

% structure: p0155 source=h2 target=section reason=relative-heading-rank; base=h2

\section{第七十六章 论闲谈之人及沉默}

当某事对某人看来无关紧要，而又不被回避时，它便轻易冲口而出，而你却滥用它。当你前来倾吐祈祷，或为每日之罪捶胸时，闲谈助长了它。号手的号声响起，而读经者正在诵读，目的是使耳朵开启，你反而阻挠了它们。你用本应叹息的嘴唇寻欢作乐。把你的胸膛向邪恶关闭，或在胸中释放它们。但因为有钱财给富人带来无耻，因此每个最自信的人都由此丧亡。同样，妇女们聚集，仿佛要进入浴池。她们拥挤，把天主的殿宇当作集市。诚然，主曾恐吓祈祷之殿。主的司铎曾下令\Quote{\InnerQuote{举心向上}}，当要祈祷时，应当保持沉默。你流利地回答，而且更不戒避应许。他为献身的人民祈求至高者，免得有人丧亡，而你却转向闲谈。你嘲笑他，或减损邻人的名誉。你无纪律地说话，仿佛天主不在场——仿佛那造作万物的主既不听也不见。\par

% structure: p0157 source=h2 target=section reason=relative-heading-rank; base=h2

\section{第七十七章 致醉汉}

对醉汉我不设限制；但我宁愿是一头野兽。你在饮酒骄傲的人面前，内心退缩，持有统治者的权能，傻瓜啊，如同在独眼巨人之中。因此你在历史中喊道：\Quote{我死时也不喝。}但愿我饮最好的东西，且心中明智。宁可援助最卑微的穷人（你还寻求什么更多可滥用的？），这样你们两人都将得到振奋。你若行这些事，就为自己熄灭了地狱。\par

% structure: p0159 source=h2 target=section reason=relative-heading-rank; base=h2

\section{第七十八章 致牧者}

你既喂养他人，且藉着不懈的喂养预备了你所能的，你做得对。但仍要照顾那不能回报你的穷人：那时你的筵席必蒙独一的天主悦纳。全能者曾特别命令这样的人要得喂养。请想，当你喂养病人时，你也是在借给至高者。在这事上，主曾愿你站立在祂面前得蒙悦纳。\par

% structure: p0161 source=h2 target=section reason=relative-heading-rank; base=h2

\section{第七十九篇。致请愿者。}

你若渴望在祈祷时蒙天垂听，就当从邪恶的隐密处折断锁链；或者，你若因怜悯穷人而行善祈祷，不要怀疑你所求的必会赐予恳求者。那时，若你毫无善行而朝拜天主，就绝不要徒然祈祷。\par

% structure: p0163 source=h2 target=section reason=relative-heading-rank; base=h2

\section{第八十篇。\Place{迦萨}人的名字。}

你们这些将要与天主-基督一同成为天上居民的人，务要持守起初，从天上察看万事。愿纯朴与温良居住在你们身上。不要无故向你们虔诚的弟兄发怒，因为你们从他那里所作的，必将承受。基督所喜悦的是：死人必要复活，是的，带着他们的身体；那些在世上被火烧死的人，当六千年完毕，世界终结时，也必复活。同时，穹苍因运行改变而变化，因为那时恶人将被天火焚烧。受造物叹息着被至高天主的怒气焚烧。那些更堪当的、出生于显赫世系的人，以及被征服的\OriginalTerm{假基督}{Antichrist}治下的贵族，要遵照天主的命令在世上再次生活一千年，他们要在奴役的轭下服侍圣徒和至高者，用脖颈背负食物。而且，当统治结束时，他们要再次受审判。那些在千年期满后轻视天主的人，将被火灭绝，那时他们自己要向山说话。墓穴和坟墓中所有的肉体都按行为复原：他们被投入地狱；他们在世上承受惩罚；惩罚向他们显示，他们从天上阅读所行的事；报应按各人的行为归于永久的暴政。我无法在一篇小论文中阐明一切；学者们的好奇心会在这其中找到我的名字。\par

对于我们的作者的第二首诗，我\Emph{毫不知情}，以下细节我感谢\Person{沙夫}博士。\par

这是一首针对犹太人和外邦人的\Emph{辩护长诗}，用粗俗的六音步诗行写成，分四十七节讨论关于天主、救赎者和人类的教义。它论及圣子和圣父的名称；在此，他很可能使人指控他为\OriginalTerm{圣父受难说异端}{Patripassian}。他转而讲述福音遭遇的障碍，警告犹太人和外邦人要弃绝他们无益的敬拜，并详细阐述他所设想的末世论。现在让我逐字引述如下：——\par

\Quote{第二首诗最有趣的部分是结尾。它比第一首诗更详尽地描述了\OriginalTerm{假基督}{Antichrist}。作者预期世界的终结将随着第七次迫害而来。\OriginalTerm{哥特人}{Goths}将征服罗马并解救基督徒；但随后\Person{尼禄}将以异教假基督的身份出现，重新征服罗马，并肆虐基督徒三年半。他反过来将被来自东方的犹太人的真假基督征服；后者在打败尼禄并焚烧罗马之后，将返回犹太地区，行虚假的奇迹，并受到犹太人的崇拜。最后基督出现，即天主自己（从作者的\Emph{唯一神论}立场来看），带着失落的十二支派[?]作为他的军队，这些支派曾居住在波斯的另一边，过着幸福简朴而有德的生活。在惊异的自然现象中，他将征服假基督及其军队，使万民归化，并夺取圣城耶路撒冷。}\par

这个双重假基督的观念在\Person{拉克坦提乌斯}的\WorkTitle{Inst. Div.}，七.16及以下中再次出现。\par

这第二首诗是由\Person{皮特拉}枢机于1852年发现的。这两首诗由\Person{E. 路德维希}编辑，\Place{莱比锡}，1877年和1878年。\par

\SourceRecord{Translated by Robert Ernest Wallis.；Ante-Nicene Fathers；Vol. 4.；Edited by Alexander Roberts, James Donaldson, and A. Cleveland Coxe.；Buffalo, NY: Christian Literature Publishing Co.,；1885.；<http://www.newadvent.org/fathers/0411.htm>.}
